\RequirePackage[orthodox]{nag}
\documentclass[11pt]{article}

%% Define the include path
\makeatletter
\providecommand*{\input@path}{}
\g@addto@macro\input@path{{include/}{../../include/}}
\makeatother

\usepackage{../../include/akazachk}

\title{Avancee en apprentissage statistique}
\author{Andres Espinosa}
\begin{document}
\pgfplotsset{compat=1.18}
\maketitle

\tableofcontents
\section{Monday September 7th}
I missed last week, so this class started in the middle of an uncertainty and evidence theory lecture.

There is a difference between a statement being supported by evidence and a statement being compatible with evidence.
If I say: "The value is likely somewhere between 3 and 4", then that supports a downstream possibility and presents a domain of claims that are compatible with that previous evidence-supported claim.

We have aleatory uncertainty and epistemic uncertainty:
\begin{itemize}
    \item \textbf{Aleatory Uncertainty:} Uncertainty due to variability in the data-generating process.
    \item \textbf{Epistemic Uncertainty: } Uncertainty due to lack of knowledge about the process.
\end{itemize}

I have gaps in my knowledge—due to missing the first class—that I must take time to amend.
But we are moving on to the next slide set so I will take this opportunity to attempt to follow along with that, and add to this section as I go.

Evidence theory has a basic formal setting.
We have a question $Q$ that has one and only one correct answer, which is unknown and we lable as $X$.
We also have $\Theta$, which is the set of all possible answers and is called the frame of discernment.
We can get evidence that gives us information about what $X$ *could* be, but it only gives us partial information about it.
So the central investigation of this course—I believe it is about evidence theory—is into how that partial information can be represented, propagated, and combined.
The separation into aleatory uncertainty—inherent to the process—and epistemic uncertainty—a failure of the modeling—seems to allow us to get more meta about not only what we know but how strong our knowledge is.


\subsection{Random Sets}
In probability theory, we have random variables.
This has an analagous representation in evidence theory, called random sets.
It takes the previously finite frame $\Theta$ and extends it to work in a concept of infinite frames, with an arbitrary frame $\Theta$ and a random closed set in $\mathbb{R}^p$.

The upper inverse of an interpretation of evidence is the set of all interpretations that is compatible.
And the lower inverse of an interpretation is the set of all interpretations that imply it.
\section{Monday September 14th}
I like how evidence theory addresses the problems with probability theory.
Everything is uncertain.
It appears that fuzzy sets and possibility theory allow us to model the degree of possibility of certain events.
\subsection{Fuzzy Sets}
A set is a collection of objects.
We have some objects and either they belong to that set, or they do not.
Fuzzy sets are a generalization of sets that allow for something to be partially in a set.
Most sets in real life are not as simple as a binary subscription.
We have gradual concepts.
Things are old and young, large and small, close to and far from; but they are not so with a specific decimal accuracy.
Instead, we have fuzzy ideas about them.

A fuzzy set represents graded membership by a function
\begin{equation}
    \tilde{A} : \Theta \to [0,1]
\end{equation}
where the value $\tilde{A}(\theta)$ represents the degree of membership of $\theta$ in the set of $A$.
A crisp set $A \subseteq \Theta$ can be identifeid with its indicator function $A(\theta) = \textbf{1}_A (\theta)$
The ordinary set can then be seen as a special case of the generalized fuzzy sets.
$F(\Theta)$ represents all fuzzy subsets of $\Theta$.

The height of a fuzzy set $\tilde{A}$ is
\begin{equation}
    hgt(\tilde{A}) = \sup_{\theta \in \Theta} \tilde{A}(\theta)
\end{equation}
If the height is 1, then the fuzzy set is \textbf{normal}.
Conceptually, that means that completely certain membership is possible to the set (we can conlude with 100\% certainty that someone is tall).

We can have upper level sets ($\alpha$-cuts) of $\tilde{A}$.
\begin{gather}
    \tilde{A}^\alpha = \{ \theta \in \Theta | \tilde{A}(\theta) \geq \alpha \} \\
    \tilde{A}^{\alpha +} = \{ \theta \in \Theta | \tilde{A}(@q) > \alpha \}
\end{gather}
If we have $\alpha = 1$, then that is the \textbf{core} of the fuzzy set (the set of all points that have full membership to the set).

How can we compare fuzzy sets?

For $\tilde{A}, \tilde{B} \in F(\Theta)$, 
\begin{equation}
    \tilde{A} \subseteq \tilde{B} \text{ iff } \forall \theta \in \Theta, \tilde{A}(\theta) \leq \tilde{B}(\theta).
\end{equation}
A relation is cutworthy if forall $\alpha$, we preserve the subset relation

We can have fuzzy intervals, numbers, quantities in $\mathbb{R}$, and we can generalize it to $\mathbb{R}^{m \times n}$ to get multidimensional fuzzy quantities and vectors.

We can also have gaussian fuzzy numbers

\begin{equation}
    \tilde{x}(x) = e^{-\frac{h}{2}(x-m)^2}
\end{equation}
where $m$ is the mode and $h$ is the precision (inverse of variance).
This can also be extended to the multivariate case.

The completement of a fuzzy set $\tilde{A}$ is
\begin{equation}
    \tilde{A}^c = 1 - \tilde{A}(\theta)
\end{equation}
There is a general negation, but in this lecture we use $1-x$.
We want to keep nice properties.

Below we have the original fuzzy intersection and union formulas
\begin{gather}
    (\tilde{A} \cap_m \tilde{B})(\theta) = \min(\tilde{A}(\theta), \tilde{B}(\theta)) \\
   (\tilde{A} \cup_m \tilde{B})(\theta) = \max(\tilde{A}(\theta), \tilde{B}(\theta)) 
\end{gather}

This has the property of being idempotent and following De Morgan laws, as well as respecting the $\alpha$-cuts (the intersection/union of the alpha cuts is the same as the alpha cut of the intersection/union).

We can also define it through product and sum
\begin{gather}
    (\tilde{A} \cap_p \tilde{B})(\theta) = \tilde{A}(\theta) \tilde{B}(\theta) \\
   (\tilde{A} \cup_p \tilde{B})(\theta) = \tilde{A}(\theta)  + \tilde{B}(\theta) - \tilde{A}(\theta) \tilde{B}(\theta)
\end{gather}

We can generalize the intersection and union operations with a $t$-norm and a $t$-conorm, as long as it respects some properties.

\subsubsection{Possibility Theory}
Fuzzy sets capture the concept of ambiguity or vagueness in claims.
Possibility theory is for us to capture the uncertainty of information.

A possibility distribution is a normal fuzzy set that expresses a fuzzy restriction on $X$
\begin{equation}
    \tilde{\pi} : \Theta \to [0,1]
\end{equation}
This measures the possibility that the candidate answer $\theta$ is the true answer $X$.

The possibility measure of a candidate set of answers $A$
\begin{equation}
    \Pi(A) = \sup_{\theta \in A} \tilde{\pi} (\theta)
\end{equation}

We can have a necessity measure that measures the extent to which $A$ is certainly true:
\begin{equation}
    N(A) = 1 - \Pi(A^c) = \inf_{\theta \notin A} [1-\tilde{\pi}(\theta)]
\end{equation}.

When we have multiple possibility distributions that constraint our true answer $X$, we can combine them to get a distribution from them.
We take the intersection, if it is possible, and then we normalize it to 1.
\begin{equation}
    (\tilde{\pi}_1 \cap_T^* \tilde{\pi}_2)(\theta) = \frac{(\tilde{\pi}_1 \cap_T \tilde{\pi}_2)(\theta)}{\sup_{\theta^\prime \in \Theta}(\tilde{\pi}_1 \cap_T^* \tilde{\pi}_2)(\theta^\prime)}
\end{equation}
Our choice of $T$ changes our function that we use for this intersection.
We can choose it based off different properties that we want (idempotence, associative, etc).
Associativity allows sequential combination of information, and idempotence allows us to capture correlated or overlapped information without reinforcing the same signal over and over again.

When we have a gaussian fuzzy vector or gaussian fuzzy number, we can combine and extend and marginalize while keeping the result also be a gaussian fuzzy vector / number.

\subsection{Exercises}
\subsubsection{Listing cuts}

Let $\tilde{A}$ be the following fuzzy subset of
$\Theta = \{a,b,c,d,e\}$:
\begin{equation}
    \tilde{A} = \{(a,0.5),(b,0.4),(c,0.7),(d,0.8),(e,1)\}.
\end{equation}

List all $\alpha$-cuts and strong $\alpha$-cuts of $\tilde{A}$.

\textbf{Solution}

\begin{equation}
    \tilde{A}^\alpha = 
    \begin{cases}
       \{ a,b,c,d,e \} & \alpha \leq 0.4 \\
       \{ a,c,d,e \} & 0.4 < \alpha \leq 0.5 \\
       \{ c,d,e \} & 0.5 < \alpha \leq 0.7 \\
       \{ d,e \} & 0.7 < \alpha \leq 0.8 \\
       \{ e \} & 0.8 < \alpha \\
    \end{cases}
\end{equation}

\subsubsection{t-norm proof}

Prove that the following inequalities hold for every $t$-norm $T$ and every
$(x,y) \in [0,1]$:
\begin{equation}
    T_D(x,y) \leq T(x,y) \leq \min(x,y),
\end{equation}
where $T_D$ is the drastic $t$-norm defined by
\begin{equation}
    T_D(x,y) =
    \begin{cases}
        \min(x,y), & \text{if } \max(x,y)=1,\\
        0, & \text{otherwise.}
    \end{cases}
\end{equation}

\textbf{Solution}
We can take the first inequality and put it in the case equation below:
\begin{gather}
    \min(x,y) \leq T(x,y) \leq \min(x,y) \quad \text{if } \max(x,y)=1 \\
    0 \leq T(x,y) \leq \min(x,y) \quad \text{otherwise.}
\end{gather}
which simplifies to
\begin{gather}
    T(x,y) = \min(x,y) \quad \text{if } \max(x,y)=1 \\
    0 \leq T(x,y) \leq \min(x,y) \quad \text{otherwise.}
\end{gather}
We can prove these separately.
Firstly, if $\max(x,y) = 1$, then that means that either $x,y = 1$.
A T-norm must follow the boundary condition property, which specifies that $T(x,1) = x$.
Clearly, the boundary condition is then equivalent to $min(x,y)$ in this case, where one of the variables is 1.

Now we must prove that 

\begin{equation}
    0 \leq T(x,y) \leq \min(x,y) \quad \text{if } \max(x,y) < 1
\end{equation}
via the monotonicity of t-norms, we get $T(x,y) \leq \min(x,y)$, and it is also non-negative.


\end{document}